% !PRINTED 338
\weberpage{338}

\begin{center}
\gesperrt{Elfter Abschnitt.}

\textbf{\Large Ideale in quadratischen K\"orpern.}

\rule{3cm}{0.4pt}
\end{center}

\begin{center}
\textbf{\S\ 90. \gesperrt{Diskriminante des quadratischen K\"orpers.}}
\end{center}

Ein quadratischer K\"orper entsteht, wenn man dem K\"orper
der rationalen Zahlen eine Quadratwurzel $\sqrt{d}$ \gesperrt{adjungiert}, hier
kann $d$ als \gesperrt{ganze Zahl ohne quadratischen Teiler} angenommen werden. Der K\"orper ist reell oder imagin\"ar, je nachdem $d$ positiv oder negativ ist. Bezeichnet man den K\"orper der
rationalen Zahlen, den absoluten Rationalit\"atsbereich, mit $\mathfrak{R}$ und
mit $\mathfrak{R}(x,x',\dots)$ den K\"orper, der durch Adjunktion von $x,x',\dots$
zu $\mathfrak{R}$ entsteht, so ist der quadratische K\"orper $\Omega$ so zu bezeichnen:
\weq{1}{\Omega=\mathfrak{R}(\sqrt{d}).}

Jede Zahl des K\"orpers $\Omega$ kann in die Form gesetzt werden:
\weq{2}{\omega=\frac{x+y\sqrt{d}}{2},}
worin $x, y$ rationale ganze oder gebrochene Zahlen sind.

Die zu $\omega$ \gesperrt{konjugierte} Zahl
\weq{3}{\omega'=\frac{x-y\sqrt{d}}{2}}
ist gleichfalls in dem K\"orper $\Omega$ enthalten, und folglich ist $\Omega$ ein
\gesperrt{Normalk\"orper.}

Damit $\omega$ eine \gesperrt{ganze Zahl} sei, ist notwendig, da\ss
\[
\omega+\omega'=x,\quad(\omega-\omega')^2=y^2 d
\]
ganze rationale Zahlen sind, und da $d$ keinen quadratischen Teiler
haben soll, so m\"ussen $x$ und $y$ ganze Zahlen sein. Damit aber
auch $\omega$ wirklich eine ganze Zahl sei, mu\ss{} auch noch die Norm
$\frac{1}{4}(x^2-y^2 d)$ eine ganze rationale Zahl sein, d.\,h. es mu\ss
\weq{4}{x^2-y^2 d\equiv 0\pmod{4}}
sein. Ist $d\equiv 2$ oder $\equiv 3\pmod 4$, so kann diese Bedingung
nur erf\"ullt sein, wenn $x$ und $y$ gerade Zahlen sind, und wenn
wir also $x, y$ durch $2x, 2y$ ersetzen, folgt:
%CONT>>
